\documentclass[11pt]{article}
\usepackage[a4paper,margin=26mm]{geometry}
\usepackage{amsmath,amssymb}
\setlength{\parindent}{0pt}
\setlength{\parskip}{7pt}
\emergencystretch=2em
\title{Lipschitz monotonicity alone does not ensure FBF convergence\\
\large Disproof of conjecture 00000008832}
\author{ziangni-sys}
\date{4 October 2026}
\begin{document}
\maketitle
Both versions of the conjecture assert unconditional convergence of
forward-backward-forward splitting under joint Lipschitz monotonicity.
Neither states existence of a zero of the inclusion. We give an actual
FBF iteration with monotone Lipschitz operators and a step satisfying
the usual strict Lipschitz step bound, but with no weak or strong limit.
The example has no zero; it does not contradict the standard convergence
theorem that assumes a solution exists.

\textbf{Actual operators and their hypotheses.}
Work in the real Hilbert space $\mathbb R$, and set
\[
 A(x)=\{0\},\qquad B(x)=1\quad(x\in\mathbb R).
\]
Both graphs are maximal monotone. To prove this directly, for any
$c\in\mathbb R$ the graph $G_c=\{(x,c):x\in\mathbb R\}$ is monotone,
because its monotonicity inner product is zero.
If a point $(p,q)$ can belong to a monotone graph containing $G_c$,
comparison with $(p-1,c)$ and $(p+1,c)$ gives
\[
 q-c\geq0,\qquad -(q-c)\geq0.
\]
Thus $q=c$, proving maximality against every monotone graph extension.

The single-valued versions of $A$ and $B$ are both constant and hence
zero-Lipschitz. In particular each is $1$-Lipschitz, so the usual
positive Lipschitz upper bound $L=1$ is valid.

\textbf{Genuine resolvent and three stages.}
For any $\gamma>0$, the resolvent equation for $A$ is
\[
 x=u+\gamma a,\qquad a\in A(u)=\{0\}.
\]
It holds exactly when $u=x$. Thus the resolvent is the actual unique map
$J_{\gamma A}=\operatorname{id}$.
The standard FBF stages are
\[
 v=x-\gamma B(x),\qquad
 y=J_{\gamma A}(v),\qquad
 x^+=y-\gamma\bigl(B(y)-B(x)\bigr).
\]
Here $v=x-\gamma$, $y=v$, and $B(y)-B(x)=0$. Consequently
\[
 x^+=x-\gamma.
\]
The backward stage is well-defined everywhere and solves its resolvent
equation exactly; no approximate certificates are assumed.

\newpage
\textbf{All iterates and absence of a limit.}
Choose $\gamma=1/2$. It satisfies $0<\gamma<1/L$ for $L=1$,
so the example does not use a violation of the usual Lipschitz bound.
For any initial point $x_0\in\mathbb R$, induction gives
\[
 x_n=x_0-n/2\qquad(n\geq0).
\]
In particular the trajectory diverges to minus infinity.
There is no finite strong limit. One direct argument, used formally,
is that convergence of $x_n$ to $a$ would imply convergence of the shifted
sequence $x_{n+1}$ to $a$ and hence of its difference to zero. But
\[
 x_{n+1}-x_n=-1/2
\]
is a nonzero constant sequence, which cannot converge to zero.

There is also no weak limit: the identity on $\mathbb R$ is an actual
continuous linear functional, so weak convergence would imply the
strong convergence just excluded. The same argument works for every
nonzero fixed step, not just $1/2$.

\textbf{The missing hypothesis.}
The actual inclusion is
\[
 0\in A(x)+B(x)=\{1\},
\]
which has no solution. Joint monotonicity and Lipschitz continuity
do not imply nonemptiness of the zero set. Thus the source's unqualified
convergence assertion is false; a valid standard FBF result needs a
solution-existence hypothesis in addition to its operator and step
assumptions. This disproof makes no separate assertion about
adaptive backtracking or the detailed joint domain of unequal step sizes.

\textbf{Formal correspondence.}
The Lean project defines actual graphs as subsets of $\mathbb R^2$
and monotonicity using the real Hilbert inner product.
\texttt{constant\_maximal} proves inclusion-maximality against every
monotone graph extension. \texttt{joint\_lipschitz} proves actual
Mathlib \texttt{LipschitzWith} bounds for both operators.

\texttt{actual\_resolvent} proves the resolvent equation and its unique
solution. The forward, backward and final forward stages are defined
explicitly; \texttt{genuine\_backward\_stage} verifies the exact equation.
\texttt{actual\_step} and \texttt{actual\_iterates} prove the update and
all iterates for every starting point and every step.
\texttt{no\_strong\_limit} proves nonconvergence, and
\texttt{no\_weak\_limit} uses the full standard functional definition
quantified over all continuous linear functionals.
\texttt{no\_zero} checks the actual inclusion. The final theorem combines
all hypotheses and the absence of any weak limit for the safe half-step.

\textbf{Reproduction.}
Use Lean 4.19.0 and public Mathlib revision
\begin{center}\small\texttt{c44e0c8ee63ca166450922a373c7409c5d26b00b}\end{center}
Run \texttt{lake build} inside \texttt{lean/}. Axiom audits are printed,
and the validation notes record the completed build and PDF inspection.
\end{document}
